% Define colours for code listings
\definecolor{keyword}{RGB}{0,0,128}
\definecolor{comment}{RGB}{0,96,0}
\definecolor{string}{RGB}{128,0,0}
\lstset{
  language=Python,
  basicstyle=\ttfamily\small,
  keywordstyle=\color{keyword}\bfseries,
  commentstyle=\color{comment}\itshape,
  stringstyle=\color{string},
  showstringspaces=false,
  frame=single,
  framerule=0.5pt,
  numbers=left,
  numberstyle=\tiny,
  numbersep=5pt
}

\chapter{Lab 4: Deploying the Customer‑Support Agent to AgentCore Runtime}
\label{chap:Lab 4: Deploying the Customer‑Support Agent to AgentCore Runtime}


\section{Introduction and Objectives}

After completing Labs 1–3, you now have a sophisticated customer‑support
agent with memory and gateway integration.  However, the agent still
runs on your local machine, which limits reliability, scalability and
observability.  Lab 4 demonstrates how to deploy the agent to
Amazon Bedrock AgentCore Runtime—a fully managed serverless environment
that delivers enterprise‑grade reliability, automatic scaling and
comprehensive monitoring.  By the end of this lab you will:

\begin{itemize}
  \item Create a runtime‑ready entrypoint using \texttt{BedrockAgentCoreApp}.
  \item Configure the deployment with the AgentCore Starter Toolkit,
        including custom JWT authentication via Cognito.
  \item Launch the runtime, monitor deployment status and invoke the
        deployed agent with a bearer token.
  \item Understand how AgentCore Observability collects traces,
        metrics and logs for debugging and performance tuning.
\end{itemize}

Compared to the previous labs, only a few additional lines of code are
required to prepare the agent for deployment.  Most of the work
happens in infrastructure configuration, which is handled by the
starter toolkit.

\section{Prerequisites}

Before proceeding ensure that you have:

\begin{itemize}
  \item Completed Lab 3, including provisioning the AgentCore Gateway
        and Cognito user pool.
  \item Python 3.12 or newer installed, along with the packages
        \texttt{strands‑agents}, \texttt{bedrock‑agentcore‑starter‑toolkit}
        and \texttt{boto3}.
  \item Docker, Finch or Podman installed and running.  The starter
        toolkit uses containers to build and push the agent image.
  \item An AWS account with permissions to create IAM roles, CodeBuild
        projects, ECR repositories and AgentCore resources.
  \item CloudWatch Transaction Search enabled in your AWS region to
        view AgentCore Observability traces.
\end{itemize}

\section{Importing Required Libraries and Initialising Memory}

As in earlier labs, begin by importing relevant modules.  For
production deployment you will also initialise the memory manager
through the AgentCore Starter Toolkit to ensure that your agent’s
memory exists in the target region.  The code below checks whether the
memory has already been created and creates it if necessary.  This step
allows the notebook to be run independently of Lab 2.

\begin{lstlisting}[language=Python]
import boto3
from lab_helpers.utils import get_ssm_parameter
from bedrock_agentcore_starter_toolkit.operations.memory.manager import MemoryManager
from bedrock_agentcore.memory.constants import StrategyType

# Establish session and region
boto_session = boto3.Session()
REGION = boto_session.region_name

# Create or retrieve the memory resource (if the user skipped Lab 2)
memory_name = "CustomerSupportMemory"
memory_manager = MemoryManager(region_name=REGION)
memory = memory_manager.get_or_create_memory(
    name=memory_name,
    strategies=[
        {
            StrategyType.USER_PREFERENCE.value: {
                "name": "CustomerPreferences",
                "description": "Captures customer preferences and behaviour",
                "namespaces": ["support/customer/{actorId}/preferences"],
            }
        },
        {
            StrategyType.SEMANTIC.value: {
                "name": "CustomerSupportSemantic",
                "description": "Stores facts from conversations",
                "namespaces": ["support/customer/{actorId}/semantic"],
            }
        },
    ],
)
memory_id = memory["id"]
\end{lstlisting}

This snippet uses the \texttt{MemoryManager} class to create the
\texttt{CustomerSupportMemory} resource with two strategies: one for
storing user preferences and another for storing semantic facts.  It
returns a memory identifier that will later be passed to the runtime
environment.

\section{Preparing the Runtime‑Ready Agent}

To deploy an agent to AgentCore Runtime you need to provide an
\emph{entrypoint} script.  This script must define an asynchronous
function decorated with \texttt{@app.entrypoint} that accepts a
payload and an optional context.  The payload is expected to contain a
\texttt{prompt} field; the context provides request headers such as
\texttt{Authorization}.  The entrypoint is where you reconstruct your
agent, using the memory configuration and gateway tools from earlier
labs.  The code below writes such a file into
\texttt{lab\_helpers/lab4\_runtime.py}.  Four marked lines show where
runtime‑specific additions occur.

\begin{lstlisting}[language=Python]
%%writefile ./lab_helpers/lab4_runtime.py
import os
from bedrock_agentcore.runtime import BedrockAgentCoreApp  #### RUNTIME - LINE 1 ####
from strands import Agent
from strands.tools.mcp import MCPClient
from mcp.client.streamable_http import streamablehttp_client
import boto3
from strands.models import BedrockModel
from lab_helpers.utils import get_ssm_parameter
from lab_helpers.lab1_strands_agent import (
    get_return_policy,
    get_product_info,
    get_technical_support,
    SYSTEM_PROMPT,
    MODEL_ID,
)
from lab_helpers.lab2_memory import ACTOR_ID, SESSION_ID
from bedrock_agentcore_starter_toolkit.operations.memory.manager import MemoryManager
from bedrock_agentcore.memory.integrations.strands.config import AgentCoreMemoryConfig, RetrievalConfig
from bedrock_agentcore.memory.integrations.strands.session_manager import AgentCoreMemorySessionManager

REGION = boto3.session.Session().region_name

# Initialise Bedrock model
model = BedrockModel(model_id=MODEL_ID)

# Retrieve memory identifier from environment (provided at launch)
memory_id = os.environ.get("MEMORY_ID")
if not memory_id:
    raise Exception("Environment variable MEMORY_ID is required")

memory_config = AgentCoreMemoryConfig(
    memory_id=memory_id,
    session_id=str(SESSION_ID),
    actor_id=ACTOR_ID,
    retrieval_config={
        "support/customer/{actorId}/semantic": RetrievalConfig(top_k=3, relevance_score=0.2),
        "support/customer/{actorId}/preferences": RetrievalConfig(top_k=3, relevance_score=0.2),
    },
)

# Initialise the AgentCore Runtime application
app = BedrockAgentCoreApp()  #### RUNTIME - LINE 2 ####

@app.entrypoint  #### RUNTIME - LINE 3 ####
async def invoke(payload, context=None):
    """Entry point for AgentCore Runtime.

    Extracts the user prompt and JWT token from the request, retrieves
    the Gateway URL from SSM, constructs an MCP client and invokes
    the agent.  Returns a plain string response.
    """
    # Extract user prompt
    user_input = payload.get("prompt", "")

    # Extract authorisation header from context (if provided)
    request_headers = context.request_headers or {}
    auth_header = request_headers.get('Authorization', '')
    print(f"Authorization header: {auth_header}")

    # Retrieve gateway details
    gateway_id = get_ssm_parameter("/app/customersupport/agentcore/gateway_id")
    gateway_client = boto3.client("bedrock-agentcore-control", region_name=REGION)
    gateway_response = gateway_client.get_gateway(gatewayIdentifier=gateway_id)
    gateway_url = gateway_response['gatewayUrl']

    # Only proceed if both gateway URL and auth token are present
    if gateway_url and auth_header:
        try:
            mcp_client = MCPClient(lambda: streamablehttp_client(
                url=gateway_url,
                headers={"Authorization": auth_header}
            ))
            with mcp_client:
                tools = [
                    get_product_info,
                    get_return_policy,
                    get_technical_support,
                ] + mcp_client.list_tools_sync()
                agent = Agent(
                    model=model,
                    tools=tools,
                    system_prompt=SYSTEM_PROMPT,
                    session_manager=AgentCoreMemorySessionManager(memory_config, REGION),
                )
                response = agent(user_input)
                # Return plain text for synchronous invocation
                return response.message["content"][0]["text"]
        except Exception as e:
            print(f"MCP client error: {str(e)}")
            return f"Error: {str(e)}"
    else:
        return "Error: Missing gateway URL or authorization header"

if __name__ == "__main__":
    app.run()  #### RUNTIME - LINE 4 ####
\end{lstlisting}

Key takeaways from this entrypoint:

\begin{itemize}
  \item The environment variable \texttt{MEMORY\_ID} is required so that the
        runtime can connect to the correct memory resource.  This variable
        will be passed by the deployment process.
  \item The Gateway URL and JWT token are retrieved at invocation time
        from SSM and the request header, respectively.  The token is
        forwarded to the Gateway to authorise the call.
  \item The code wraps the MCP client in a context manager to ensure
        resources are cleaned up and tools are automatically discovered.
  \item The function returns plain text rather than a structured object.
        AgentCore Runtime expects the entrypoint to return a string or
        dictionary representing the response body.
\end{itemize}

\section{Configuring the AgentCore Runtime Deployment}

The AgentCore Starter Toolkit provides a high‑level API for configuring
and launching an AgentCore Runtime.  This configuration step
specifies the entrypoint file, the IAM role that the runtime will
assume, whether to automatically create an Amazon Elastic Container
Registry (ECR) repository, and the requirements file to be installed in
the container.  It also sets up JWT‑based authentication by
referencing the Cognito client ID and discovery URL stored in SSM.

\begin{lstlisting}[language=Python]
from bedrock_agentcore_starter_toolkit import Runtime
from lab_helpers.utils import create_agentcore_runtime_execution_role, get_or_create_cognito_pool

# Create an execution role granting the runtime permission to
# access AWS services (S3, DynamoDB, Bedrock, etc.)
execution_role_arn = create_agentcore_runtime_execution_role()

# Initialise the runtime toolkit
agentcore_runtime = Runtime()

response = agentcore_runtime.configure(
    entrypoint="lab_helpers/lab4_runtime.py",
    execution_role=execution_role_arn,
    auto_create_ecr=True,
    requirements_file="requirements.txt",
    region=REGION,
    agent_name="customer_support_agent",
    authorizer_configuration={
        "customJWTAuthorizer": {
            "allowedClients": [get_ssm_parameter("/app/customersupport/agentcore/client_id")],
            "discoveryUrl": get_ssm_parameter("/app/customersupport/agentcore/cognito_discovery_url"),
        }
    },
    # Forward specific headers to the runtime to propagate OAuth tokens
    request_header_configuration={
        "requestHeaderAllowlist": [
            "Authorization",
            "X-Amzn-Bedrock-AgentCore-Runtime-Custom-H1",
        ]
    },
)
print("Configuration completed:", response)
\end{lstlisting}

This call generates two artefacts: a \texttt{Dockerfile} that defines
how the container is built and a \texttt{.bedrock\_agentcore.yaml}
configuration file.  The IAM execution role provides runtime access to
AWS services, while the custom JWT authorizer enforces that only
clients with valid Cognito tokens can invoke the agent.  The header
allowlist ensures that the \texttt{Authorization} header is passed from
the caller to the entrypoint so that tokens are not stripped out by
API Gateway or other intermediaries.

\section{Launching the Runtime and Checking Status}

After configuration, call \texttt{launch} to build and deploy the
container image.  Behind the scenes the starter toolkit uses AWS
CodeBuild to build the image, pushes it to ECR and creates the
AgentCore Runtime resources.  Pass the memory identifier via
\texttt{env\_vars} so that the entrypoint can find the correct
memory.

\begin{lstlisting}[language=Python]
from lab_helpers.utils import put_ssm_parameter

# Launch the agent (builds and deploys the container)
launch_result = agentcore_runtime.launch(env_vars={"MEMORY_ID": memory_id})
print("Launch completed:", launch_result.agent_arn)

# Persist the runtime ARN for later use
put_ssm_parameter("/app/customersupport/agentcore/runtime_arn", launch_result.agent_arn)

# Wait for the runtime to become READY
import time
status_response = agentcore_runtime.status()
status = status_response.endpoint["status"]
end_status = ["READY", "CREATE_FAILED", "DELETE_FAILED", "UPDATE_FAILED"]
while status not in end_status:
    print(f"Waiting for deployment... Current status: {status}")
    time.sleep(10)
    status_response = agentcore_runtime.status()
    status = status_response.endpoint["status"]
print(f"Final status: {status}")
\end{lstlisting}

The \texttt{status} method polls the runtime’s state until it reaches
\texttt{READY} or a failure condition.  If the runtime already
exists, you can pass \texttt{auto\_update\_on\_conflict=True} to
\texttt{launch} to update the existing deployment.  Once the status is
\texttt{READY}, the runtime exposes an invocation endpoint that can
handle concurrent requests at scale.

\section{Invoking the Deployed Agent}

With the runtime active, you can test it by sending requests via the
starter toolkit.  You must supply a valid Cognito access token and a
session identifier.  The toolkit will automatically propagate the
token as an Authorization header and create a unique session ID if
none is provided.  The examples below demonstrate both session
continuity and new sessions.

\begin{lstlisting}[language=Python]
import uuid

# Retrieve an access token (must still be valid)
from lab_helpers.utils import get_or_create_cognito_pool
access_token = get_or_create_cognito_pool(refresh_token=True)

# Extract runtime ID from the ARN
runtime_id = launch_result.agent_arn.split(":")[-1].split("/")[-1]
print(f"Runtime ID: {runtime_id}")

session_id = str(uuid.uuid4())

# First invocation: list tools
prompt = "List all of your tools"
response = agentcore_runtime.invoke(
    {"prompt": prompt},
    bearer_token=access_token["bearer_token"],
    session_id=session_id,
)
print(response["response"])

# Second invocation with the same session (session continuity)
prompt2 = "Tell me detailed information about the technical documentation on installing a new CPU"
response = agentcore_runtime.invoke(
    {"prompt": prompt2},
    bearer_token=access_token["bearer_token"],
    session_id=session_id,
)
print(response["response"])

# Third invocation with a new session (new customer)
session_id2 = str(uuid.uuid4())
prompt3 = "I have a Gaming Console Pro device, I want to check my warranty status, warranty serial number is MNO33333333."
response = agentcore_runtime.invoke(
    {"prompt": prompt3},
    bearer_token=access_token["bearer_token"],
    session_id=session_id2,
)
print(response["response"])
\end{lstlisting}

The first call lists all available tools (three local and two MCP) to
verify that the runtime is correctly configured.  The second call uses
the same session ID, allowing the agent to remember previous context
and memory state.  The third call uses a different session ID,
simulating a new customer; the agent will not have prior context and
will ask clarifying questions if needed.

\section{Observability with CloudWatch}

AgentCore Runtime integrates out‑of‑the‑box with AgentCore
Observability to send traces, metrics and logs to CloudWatch.  Once
your runtime is running, you can explore these observability features
in the AWS console under \texttt{CloudWatch \rightarrow GenAI
Observability \rightarrow Bedrock AgentCore}.  There you will find
dashboards for:

\begin{itemize}
  \item \textbf{Agents}: a high‑level overview showing throughput,
        latency and error rates across your deployed agents.
  \item \textbf{Sessions}: a list of session IDs with details on
        conversation length, tool usage and memory retrievals.
  \item \textbf{Traces}: deep traces of individual requests showing
        model invocation, tool calls, memory operations and latency
        breakdown.  You can filter, search and export traces for
        debugging and performance tuning.
\end{itemize}

These dashboards help you identify performance bottlenecks, monitor
tool invocation patterns and ensure that your agent behaves as
expected in production.  You can also query logs using CloudWatch
Logs Insights or integrate with other observability tools if needed.

\section{Summary and Next Steps}

In this lab you transformed your agent into a production‑ready service
running on Amazon Bedrock AgentCore Runtime.  You wrote an entrypoint
script using \texttt{BedrockAgentCoreApp}, configured the runtime
deployment with the starter toolkit, launched the container, tested
session continuity and observed how to monitor the agent using
CloudWatch.  The runtime automatically scales to handle concurrent
requests, provides secure access via JWT tokens and offers
comprehensive observability.

Remaining challenges include providing a user‑friendly interface for
end customers and evaluating the agent’s performance over time.
These topics are addressed in Labs 5 and 6.  In the next lab you will
configure online evaluations to continuously monitor quality metrics
and learn how to interpret the results.

